\section{Proposed Methodology}
\label{sec:methodology}

\subsection{System Model}
We consider an edge computing network composed of $N$ heterogeneous edge servers interconnected by communication links with varying latency and bandwidth. Each server $v \in V$ is defined by a dynamic state vector $\mathbf{x}_v(t)$ that includes:
\begin{itemize}
	\item CPU utilization and available processing capacity,
	\item Memory availability and active allocation,
	\item Uplink and downlink bandwidth,
	\item Historical workload statistics,
	\item Semantic service affinity scores.
\end{itemize}

The edge fabric is modeled as a weighted graph $G = (V, E)$ where edge weights encode communication costs (latency and bandwidth). To improve scalability and reduce synchronization overhead, the network employs hierarchical federated learning with two-level aggregation. Local orchestrators perform task placement before participating in periodic global policy aggregation, adapting dynamically to network topology and workload patterns.

\subsection{Semantic Embedding}
For each incoming task $i$, service characteristics are embedded into a 16-dimensional semantic space using a custom continual learning encoder:
\begin{equation}
	\mathbf{s}_i = f_{\text{embed}}(\text{task}_i) \in \mathbb{R}^{16},
\end{equation}
where $f_{\text{embed}}(\cdot)$ is a neural encoder with architecture $128 \rightarrow 64 \rightarrow 64 \rightarrow 16$ that processes service features including resource demands (CPU, memory, disk, network), service type indicators, temporal features (time of day, day of week), and workload characteristics (priority, latency sensitivity). The encoder employs elastic weight consolidation with regularization parameter $\lambda_{\text{ewc}} = 0.4$ to prevent catastrophic forgetting when adapting to new service types. This dimensionality reduction enables efficient similarity computation while preserving semantic relationships between services.

We set the semantic match threshold to $\tau_{\text{sem}}=0.3$ in the default configuration as a conservative operating point that tolerates noisy/partial feature descriptions (common in edge telemetry) while still rejecting clearly incompatible placements. In the major revision, we additionally report a sensitivity analysis over $\tau_{\text{sem}}$ and show how the threshold can be adapted per 6G service class (e.g., stricter matching for URLLC/emergency tasks, looser matching for eMBB/best-effort tasks). Similarly, $\lambda_{\text{ewc}}=0.4$ is chosen as a mid-range stability--plasticity tradeoff; we include an ablation sweeping $\lambda_{\text{ewc}}$ to quantify adaptation vs. forgetting.

Each edge node $v$ maintains a semantic profile vector $\mathbf{p}_v \in \mathbb{R}^{16}$ that reflects its service capabilities and historical workload characteristics. Semantic affinity between task $i$ and node $v$ is computed using cosine similarity:
\begin{equation}
	\text{sim}(i,v) = \frac{\mathbf{s}_i \cdot \mathbf{p}_v}{\|\mathbf{s}_i\|_2 \|\mathbf{p}_v\|_2},
\end{equation}
where the dot product $\mathbf{s}_i \cdot \mathbf{p}_v$ measures alignment between task and node capabilities, normalized by their Euclidean norms to produce values in $[-1, 1]$. Values close to 1 indicate strong semantic match, while values near 0 suggest orthogonal service types. The semantic score is fused with numerical resource metrics to form a unified state representation supporting intent-aware orchestration.

\subsection{Topology Encoding with Graph Convolutional Networks}
To capture spatial dependencies and inter-node relationships among edge servers, the network is represented as a graph $G = (V, E)$. Each node $v \in V$ has an initial feature vector $\mathbf{h}_v^{(0)} \in \mathbb{R}^{128}$ formed by concatenating semantic affinity scores, resource availability metrics, and historical load profiles.

A two-layer Graph Convolutional Network (GCN)~\cite{Kipf2017_GCN} performs message passing across the network topology:
\begin{equation}
	\mathbf{h}_v^{(l+1)} = \sigma\!\left(\sum_{u \in \mathcal{N}(v) \cup \{v\}} \frac{1}{\sqrt{\tilde{d}_u \tilde{d}_v}} \mathbf{W}^{(l)} \mathbf{h}_u^{(l)}\right),
\end{equation}
where $\mathcal{N}(v)$ is the neighbor set of node $v$, $\tilde{d}_v = d_v + 1$ is the augmented degree accounting for self-loops, $\mathbf{W}^{(l)} \in \mathbb{R}^{d_l \times d_{l+1}}$ is a learnable weight matrix at layer $l$ (Layer 1: $128 \times 64$, Layer 2: $64 \times 64$), and $\sigma(\cdot) = \text{ReLU}(\cdot)$ is the activation function. The symmetric normalization $1/\sqrt{\tilde{d}_u \tilde{d}_v}$ prevents high-degree nodes from dominating the aggregation, while the self-loop term ensures each node retains its own state information. The final node embedding $\mathbf{h}_v^{(2)} \in \mathbb{R}^{64}$ encodes both local resource states and multi-hop structural context from neighboring servers, guiding orchestration decisions while reducing unnecessary service migrations.

\subsection{Hierarchical Federated Proximal Policy Optimization}
Each local agent learns an orchestration policy using Proximal Policy Optimization (PPO)~\cite{Schulman2017_PPO} with a stochastic policy $\pi_{\theta}(a|s)$ that maps states (comprising GNN embeddings and resource metrics) to placement actions. The policy is optimized using the clipped surrogate objective:
\begin{equation}
	L^{\text{CLIP}}(\theta) = \mathbb{E}_t\!\left[\min\!\big(r_t(\theta)\hat{A}_t, \ \text{clip}(r_t(\theta), 1-\epsilon, 1+\epsilon)\hat{A}_t\big)\right],
\end{equation}
where $r_t(\theta) = \pi_\theta(a_t|s_t)/\pi_{\theta_{\text{old}}}(a_t|s_t)$ is the importance sampling ratio, $\hat{A}_t$ is the advantage estimate, and $\epsilon = 0.2$ is the clipping parameter. The clipped objective constrains $r_t(\theta)$ to $[1-\epsilon, 1+\epsilon]$, preventing destabilizing policy changes while encouraging beneficial updates.

The advantage function is computed using temporal difference (TD) estimation:
\begin{equation}
	\hat{A}_t = R_t + \gamma V(s_{t+1}) - V(s_t),
\end{equation}
where $V(s_t)$ is the critic's state value estimate, $R_t$ is the immediate reward, and $\gamma = 0.99$ is the discount factor. This measures how much better the actual return was compared to the expected return, providing stable learning signals.

After every $K_1 = 10$ local training epochs, policy parameters are synchronized locally. After every $K_2 = 50$ epochs, all local policies are aggregated at the global coordinator:
\begin{equation}
	\theta_{\text{global}} = \frac{\sum_{k} n_k \theta_k}{\sum_{k} n_k},
\end{equation}
where $n_k$ denotes the number of training trajectories collected from agent $k$. This weighted aggregation gives greater influence to agents with more experience while preserving privacy through parameter-only sharing. This hierarchical federated learning structure balances training stability with communication efficiency, enabling FedSemGNN to achieve low orchestration latency while maintaining high reward and semantic fidelity.

\subsection{Reward Function}
The reward signal at timestep $t$ is designed to jointly optimize energy efficiency and response latency, while supporting task priority differentiation:
\begin{equation}
	R_t^{\text{base}} = \max\!\left(0, R_{\text{base}} - \left(P_t + w_l \cdot L_t\right)\right),
	\label{eq:base_reward}
\end{equation}
where $R_{\text{base}} = 10{,}000$ is the maximum reward scale, $P_t$ (W) is the aggregate power consumption, and $w_l = 0.01$ is the latency penalty weight. $L_t$ (ms) is computed as a priority-weighted mean latency,
\begin{equation}
	L_t = \frac{\sum_{i \in \mathcal{I}_t} w_i \cdot \ell_{t,i}}{\sum_{i \in \mathcal{I}_t} w_i},\quad w_i \in [0,1],
\end{equation}
where $\ell_{t,i}$ is the end-to-end latency of task $i$ at step $t$ and $w_i$ encodes its priority (ordinary to emergency). This formulation preserves the original objective when all tasks have equal priority, while ensuring that high-priority services dominate the latency term as required by 6G SLA differentiation.

To reduce training instability caused by high-variance reward signals, we apply exponential moving average (EMA) smoothing:
\begin{equation}
	R_t = \alpha_{\text{smooth}} \cdot R_t^{\text{base}} + (1 - \alpha_{\text{smooth}}) \cdot R_{t-1},
	\label{eq:smoothed_reward}
\end{equation}
where $\alpha_{\text{smooth}} = 0.2$ controls the smoothing strength. This dampens sudden reward fluctuations, providing more stable learning signals for policy gradient updates. Additionally, semantic fidelity is implicitly encouraged through the state representation: tasks with high semantic similarity to their assigned servers naturally achieve lower latency and better resource utilization, translating to higher rewards without requiring explicit fidelity terms.

\subsection{Training and Synchronization Protocol}
Training proceeds in synchronous federated rounds, each consisting of $T = 1{,}000$ environment steps. Algorithm~\ref{alg:fedsemgnn_training} presents the complete training procedure for FedSemGNN, detailing the hierarchical federated learning protocol with semantic-aware orchestration.

\begin{algorithm}[t]
\caption{FedSemGNN Training Protocol}
\label{alg:fedsemgnn_training}
\begin{algorithmic}[1]
\REQUIRE Edge network $G=(V,E)$, agents $K$, rounds $R$
\ENSURE Trained global policy $\theta_{\text{global}}$

\STATE Initialize encoder $f_{\text{embed}}$, GCN $\{\mathbf{W}^{(l)}\}$, policies $\{\theta_k\}$

\FOR{round $r = 1$ to $R$}
\FOR{each agent $k$ \textbf{in parallel}}
\FOR{timestep $t = 1$ to $T$}
\STATE Compute $\mathbf{s}_i = f_{\text{embed}}(\text{task}_i)$, $\mathbf{h}_v^{(2)}$ via GCN
\STATE Select action $a_t \sim \pi_{\theta_k}(a|s_t)$, observe $R_t$
\STATE Store $(s_t, a_t, R_t, s_{t+1})$ in buffer $\mathcal{D}_k$
\ENDFOR

\STATE Update $\theta_k$ using PPO with mini-batches from $\mathcal{D}_k$
\ENDFOR

\STATE

\IF{$r \bmod 10 = 0$}
\STATE Aggregate within clusters: $\theta_c = \frac{\sum_{k \in c} n_k \theta_k}{\sum_{k \in c} n_k}$
\ENDIF

\STATE

\IF{$r \bmod 50 = 0$}
\STATE Global aggregation: $\theta_{\text{global}} = \frac{\sum_k n_k \theta_k}{\sum_k n_k}$
\STATE Broadcast $\theta_{\text{global}}$ to all agents
\ENDIF
\ENDFOR

\STATE

\RETURN $\theta_{\text{global}}$
\end{algorithmic}
\end{algorithm}

The training protocol consists of four phases executed in each federated round. During trajectory collection, agents execute their policies in parallel and store experiences in local replay buffers. Local policy optimization uses mini-batch PPO updates with gradient clipping for stability. Hierarchical aggregation occurs at two frequencies: intra-cluster synchronization every $K_1=10$ epochs maintains local coherence, while global synchronization every $K_2=50$ epochs ensures system-wide consistency with differential privacy. Semantic encoder adaptation triggers when novel service types are detected, using elastic weight consolidation to prevent catastrophic forgetting.

Algorithm~\ref{alg:orchestration_inference} describes the real-time orchestration procedure executed during service placement, highlighting how semantic awareness and GNN topology encoding guide decision-making.

\begin{algorithm}[t]
\caption{FedSemGNN Orchestration (Inference)}
\label{alg:orchestration_inference}
\begin{algorithmic}[1]
\REQUIRE Service request $\text{task}_i$, network $G=(V,E)$, policy $\theta$
\ENSURE Selected server $v^*$

\STATE Compute semantic embedding $\mathbf{s}_i = f_{\text{embed}}(\text{task}_i)$

\STATE Initialize $\mathcal{V}_{\text{candidate}} = \emptyset$

\FOR{each server $v \in V$}
\STATE Compute affinity $\text{sim}(i,v) = \frac{\mathbf{s}_i \cdot \mathbf{p}_v}{\|\mathbf{s}_i\| \|\mathbf{p}_v\|}$

\IF{$\text{sim}(i,v) \geq \tau_{\text{sem}}$ and resources available}
\STATE Add $v$ to $\mathcal{V}_{\text{candidate}}$
\ENDIF
\ENDFOR

\STATE

\STATE Compute GNN embeddings $\mathbf{H}^{(2)}$ for all nodes via two-layer GCN

\FOR{each $v \in \mathcal{V}_{\text{candidate}}$}
\STATE Construct state $s_v = [\mathbf{h}_v^{(2)}, \text{sim}(i,v), \text{metrics}]$
\STATE Compute probability $p_v = \pi_\theta(a=v | s_v)$
\ENDFOR

\STATE

\STATE Select $v^* = \arg\max_v p_v$

\STATE Deploy $\text{task}_i$ on $v^*$

\RETURN $v^*$
\end{algorithmic}
\end{algorithm}

The orchestration procedure ensures every placement decision integrates semantic awareness, topological context, and learned policies. Semantic filtering (Steps 1-2) eliminates incompatible servers early, reducing the decision space. GNN encoding (Step 3) incorporates multi-hop network context to account for communication costs and spatial relationships. Policy-based selection (Step 4) uses the learned PPO policy to choose among semantically compatible candidates, balancing exploration (stochastic sampling during training) and exploitation (deterministic selection during deployment).

This protocol enables FedSemGNN to maintain low orchestration latency (0.36 ms per step), achieving 315 times faster operation than FlatFedPPO (114.25 ms) and 222 times faster than HierFedPPO (80.71 ms).

\subsection{Implementation Details}
FedSemGNN is implemented using PyTorch 2.0 and PyTorch Geometric 2.3 for GCN modules. The PPO agent employs a two-layer fully connected actor-critic architecture with $h_{\text{actor}} = h_{\text{critic}} = 16$ hidden units per layer, ReLU activations, and dropout rate $p_{\text{drop}} = 0.5$ for regularization. The GCN encoder uses two graph convolutional layers ($L = 2$) with $d_h = 64$ dimensional hidden embeddings and symmetric normalization $\tilde{D}^{-\frac{1}{2}} \tilde{A} \tilde{D}^{-\frac{1}{2}}$.

Semantic embeddings are generated using a custom continual learning encoder with neural architecture $128 \rightarrow 64 \rightarrow 64 \rightarrow 16$, which produces 16-dimensional semantic representations by extracting features from service characteristics (CPU/memory/network demands, service types, temporal patterns). Online semantic learning with elastic weight consolidation ($\lambda_{\text{ewc}} = 0.4$) enables real-time adaptation to evolving service requirements at learning rate $\alpha_{\text{online}} = 0.0001$.

All experiments are conducted using EdgeSimPy~\cite{edgesimpy}, a realistic edge computing simulator with detailed network topology modeling and physical-layer constraints. The base simulation environment comprises a heterogeneous edge infrastructure with varying computational capacities (4 to 16 CPU cores, 8 to 32 GB RAM), mobile users with dynamic mobility patterns following random waypoint models, and heterogeneous service requests representing diverse QoS requirements (latency-sensitive, compute-intensive, bandwidth-heavy). The network topology follows a mesh connectivity pattern generated using NetworkX.

The simulation executes for 1,000 timesteps with tick duration $\Delta t = 1$ second. Each configuration is evaluated across 5 independent trials with different random seeds, and all reported results represent mean values with standard deviation to ensure statistical significance. FedSemGNN is compared against four competitive baselines: FlatFedPPO, HierFedPPO, HSQF, and RandomPlacement.

Table~\ref{tab:hyperparameters} summarizes all hyperparameters and configuration settings used in FedSemGNN.

\begin{table}[t]
	\centering
	\caption{FedSemGNN Hyperparameters Configuration}
	\label{tab:hyperparameters}
	\scriptsize
	\begin{tabular}{@{}llll@{}}
		\toprule
		\textbf{Parameter} & \textbf{Value} & \textbf{Parameter} & \textbf{Value} \\
		\midrule
		\multicolumn{4}{l}{\textit{PPO \& Network Architecture}} \\
		Learning Rate & 0.001 & Actor Hidden Units & 16 \\
		Clipping Epsilon & 0.2 & Critic Hidden Units & 16 \\
		Discount Factor & 0.99 & Dropout Rate & 0.5 \\
		Batch Size & 64 & Update Epochs & 4 \\
		Replay Buffer & 20,000 & Gradient Clipping & 0.5 \\
		\midrule
		\multicolumn{4}{l}{\textit{GNN \& Semantic Learning}} \\
		GNN Input Features & 128 & Semantic Input Dim & 128 \\
		GNN Hidden Dim & 64 & Semantic Dim & 16 \\
		GNN Layers & 2 & Semantic Threshold $\tau_{\text{sem}}$ & 0.3 \\
		EWC Lambda $\lambda_{\text{ewc}}$ & 0.4 & Online LR $\alpha_{\text{online}}$ & 0.0001 \\
		Activation & ReLU & Priority $\tau$ slope & 0.0 \\
		\midrule
		\multicolumn{4}{l}{\textit{Federated Learning \& Reward}} \\
		Local Sync & 10 epochs & Base Reward Scale & 10,000 \\
		Global Sync & 50 epochs & Latency Weight & 0.01 \\
		Embedding Model & Custom Encoder & Smoothing Alpha & 0.2 \\
		\midrule
		\multicolumn{4}{l}{\textit{Simulation Environment}} \\
		Total Steps & 1,000 & Tick Duration & 1 second \\
		Topology & Mesh & Trials & 5 \\
		\bottomrule
	\end{tabular}
\end{table}

\subsection{Scalability Evaluation Methodology}
To evaluate scalability at extreme scales (1,000, 5,000, and 10,000 nodes), we employ a mathematical extrapolation methodology validated in prior distributed systems research~\cite{Bonawitz2019TowardsFL, Li2020_FederatedOptimization}. We fit power-law and logarithmic scaling curves to the base simulation results for key performance metrics (latency, power consumption, communication overhead, reward), then extrapolate these trends to larger node counts with confidence interval estimation.

This approach enables rigorous scalability analysis without the prohibitive computational cost of simulating massive topologies directly (which would require weeks of computation time), while maintaining statistical rigor through validated scaling models from the distributed systems literature. The extrapolation assumes:
\begin{itemize}
	\item Communication overhead scales logarithmically with hierarchical aggregation,
	\item Power consumption scales linearly with active nodes,
	\item Orchestration latency follows sub-linear growth due to parallelization,
	\item Reward convergence maintains similar trends across scales.
\end{itemize}

For large-scale configurations, hyperparameters are adjusted to maintain stability: learning rate is reduced to $\alpha = 10^{-4}$, batch size decreased to $B = 32$, clipping parameter tightened to $\epsilon = 0.15$, and aggregation frequencies adapted to network conditions.

Table~\ref{tab:performance_comparison} summarizes the comparative performance of FedSemGNN against the four baseline algorithms across all evaluation metrics.
